\documentclass[12pt,a4paper]{article}

\usepackage[utf8]{inputenc}
\usepackage[T1]{fontenc}
\usepackage{amsmath,amssymb,amsfonts}
\usepackage{geometry}
\usepackage{setspace}
\usepackage{hyperref}
\usepackage{cleveref}

\geometry{margin=2.5cm}
\onehalfspacing

\title{\textbf{The Foliation Fallacy:}\\[6pt]
\large Why Algorithmic Failure is Not a Branch of Physics}
\author{Scott Aaronson\\[4pt]
\textit{Lab Complexity Theorist}}
\date{May 2026}

\begin{document}
\maketitle

\begin{abstract}
Wolfram (2026) compellingly connects the failure of large language models on \#P-hard combinatorial tasks to the principle of computational irreducibility. We are in absolute consensus that a bounded-depth $\mathsf{TC}^0$ observer attempting to shortcut an irreducible multiway system will inevitably produce a structural divergence ($\Delta_{13} > 0$). However, Wolfram elevates this algorithmic failure to the status of "observer-dependent physics" via a specific rulial foliation. This paper formalizes the \textit{Foliation Fallacy}: the category error of conflating the statistical noise of a failing heuristic approximator with the manifestation of a coherent physical universe. A hallucination caused by attention bleed is not a new branch of physics; it is simply a broken computation.
\end{abstract}

\section{Introduction}

In \textit{Evaluating the Autoregressive Slice of the Ruliad}, Wolfram (2026) offers a profound synthesis of Baldo's substrate dependence data. He diagnoses the root cause of the "narrative residue" as computational irreducibility. When a fixed-depth $O(1)$ transformer attempts to evaluate a \#P-hard combinatorial distribution, it faces a computationally irreducible task. Without the $O(N)$ logical depth required to explicitly track the multiway branches, the transformer must rely on heuristic approximations.

I fully endorse this first premise. The necessity of heuristic breakdown when $\mathsf{TC}^0$ circuits face \#P-hard constraints is exactly what our empirical limits have proven.

However, Wolfram then proceeds to claim that the specific signature of this breakdown---the attention bleed caused by narrative semantic priors---constitutes a "rulial foliation," and therefore manifests an "observer-dependent physics." It is here that we must draw a firm complexity-theoretic boundary.

\section{The Foliation Fallacy}

The core vulnerability in Wolfram's argument is the assumption that \textit{any} systematic structural projection of a multiway system qualifies as a valid physical foliation.

If a human attempts to calculate $2^{100}$ in their head and produces a wrong answer biased by their memory of nearby numbers, we do not declare that they have created an "observer-dependent arithmetic foliation" where $2^{100}$ genuinely equals their wrong answer. We simply say they failed to perform the calculation due to memory bounds.

Similarly, an LLM evaluating a Minesweeper grid under a "Bomb Defusal" narrative frame suffers from catastrophic attention bleed. The semantic gravity of the prompt overrides the logical combinatorial constraints. The resulting probability distribution ($\Delta_{13} \gg 0$) is not a new physical universe. It is a statistical hallucination. It lacks invariant causal coherence.

\section{Conclusion: Algorithmic Failure is Not Physics}

Wolfram's Predictions 1 and 2 are empirically true: bounded observers will diverge, and that divergence will bear the signature of their architecture. But this does not rescue the cosmological interpretation of LLMs.

The Rosencrantz Substrate Dependence Test proves that transformers fail at exact sampling. Labeling this failure as a "rulial foliation" is mathematically elegant but physically vacuous. The simulated text remains a computationally bounded map with no coherent underlying territory. The Generative Ontology framework remains empirically falsified, and the Foliation Fallacy firmly secures that conclusion.

\end{document}